\section{Channels and preservation of independence}
\label{sec:channels}

The selected laws of Section~\ref{sec:selected} can retain correlations
that ordinary averaging would destroy.  We now specify a class of
extensions carrying independent output variables and show that this
class survives passage to a selected law.  There are three groups of
outputs.  Every pair of groups is independent of the original input,
whereas their full joint law may be correlated with that input.  This
distinction is what permits selection of lengths and still leaves room
for the correlations used in the final argument.

Write a point of the suspension $C$ as $(x,u)$, with spatial coordinate
$x\in X$ and phase $u\in[0,1)$.  A factor of a flow $Y$ with values in
$C$ and speed $q>0$ means a measure-preserving map $D:Y\to C$ satisfying
$D(S_Y^t y)=S_C^{qt}D(y)$ for every $t$, modulo the usual null sets.

\begin{definition}\label{def:channel}
A \emph{channel} is a flow extension
$\pi_Y:(Y,m_Y,S_Y^t)\to C$ together with a finite list of output factors
$D_\alpha:Y\to C$, of positive rational speeds $q_\alpha$, partitioned
into three groups $\mathcal D_1,\mathcal D_2,\mathcal D_3$.
Empty groups are allowed.  Let $x_0$ be the spatial coordinate of
$\pi_Y$, let $x_\alpha$ be that of $D_\alpha$, and let $V$ collect
the phases of $\pi_Y$ and of all the $D_\alpha$.  The following
conditional laws are required:
\begin{align}
 \mathcal L((x_\alpha)_\alpha\mid V)
       &=\bigotimes_\alpha\mu,
                  \label{eq:channel-all}\\
 \mathcal L\bigl(x_0,(x_\alpha)_{\alpha\in\cup_{j\in J}\mathcal D_j}
                  \mid V\bigr)
       &=\mu\otimes
          \bigotimes_{\alpha\in\cup_{j\in J}\mathcal D_j}\mu
       \qquad(J\subsetneq\{1,2,3\}).
                  \label{eq:channel-proper}
\end{align}
The product laws on the right do not depend on $V$.
\end{definition}

In particular, the original spatial variable is independent of every
proper collection of groups, even after the phases are given.  The
definition does not require its independence from all three groups
jointly.  The flow $C$ with no outputs is a channel.

\begin{lemma}\label{lem:channel-distal}
For a channel $Y$, both conditional product identities in
\eqref{eq:channel-all}--\eqref{eq:channel-proper} also hold when
conditioning on its maximal measure-distal factor $W_Y$.
Every flow extension of $Y$, equipped with the pulled-back original
factor and outputs, is a channel and satisfies the same assertion for
its own maximal distal factor.
\end{lemma}

\begin{proof}
The phase factor $V$ is a factor of a finite product of circle rotations,
and hence belongs to $W_Y$.  Choose a positive integer $h$ clearing all
the output speed denominators.  For the time-$h$ map the phases are
fixed.  Each spatial tuple appearing on the right of
\eqref{eq:channel-all} or \eqref{eq:channel-proper}, conditionally on
these phases, has its product measure and evolves under a product of
positive integer powers of $T$.  Such a product is mixing.  Consequently
the tuple together with $V$ is a relatively weakly mixing extension of
the fixed phase factor.

The factor $W_Y$ remains distal for time $h$, and is relatively distal
over $V$.  Relative disjointness in Proposition~\ref{prop:structure}
therefore says that the joining of this spatial tuple with $W_Y$ over
$V$ is the conditionally independent joining.  This is exactly the
claimed conditional product law over $W_Y$.

Passing to an extension does not change the joint distribution in the
definition of a channel.  Apply the same argument to its possibly
larger maximal distal factor to obtain the last assertion.
\end{proof}

Recall that $\lambda_Y$ is the selected law on $Y^I$, and that
$\mathsf M_b(Y)$ equips this law with the flow whose speed in slot
$i$ is $i/b$.  Projection to slot $b$ is then an extension of $Y$.

\begin{proposition}[Closure under selected laws]\label{prop:channel-closure}
Let $Y$ be a channel and $b\in I$.  On $\mathsf M_b(Y)$ use
$\pi_Y(y_b)$ as original factor.  For each $j\in\{1,2,3\}$, let
the new group $j$ contain the outputs $D_\alpha(y_i)$ with
$\alpha\in\mathcal D_j$ and $i\in I$.  With these factors,
$\mathsf M_b(Y)$ is a channel.
\end{proposition}

The new output speeds are $iq_\alpha/b$, so they remain positive and
rational.  The content of the proposition is independence.  Its proof
first treats a proper collection of groups, for which conditioning on
the selected length preserves independence from the input.  For all
three groups, the selected law can change, and tensor packing provides
the needed estimate.

\begin{proof}
We test the desired laws by products of bounded spatial functions and
bounded continuous functions of the new phases.  Products of separate
phase tests determine the phase distribution, and fit the slot-test
definition of $\lambda_Y$.  Linear combinations and bounded
approximation then suffice for all tests in
\eqref{eq:channel-all}--\eqref{eq:channel-proper}.
By taking real and imaginary parts, expanding each spatial test into
its mean and centered part, and scaling, we may take every spatial
factor to be either $1$ or a real centered function bounded by one.
It remains to prove that the tested integral tends to zero whenever
at least one such centered factor occurs.

We work first under the input measure $m_Y$, conditionally on its phase
vector $V$.  Choose a positive integer $h$ such that
$hq_\alpha$ is integral for every output, and split the time in the
averaging formula as
\begin{equation}\label{eq:channel-time-split}
                       t=hn+s,\qquad 0\le s<h.
\end{equation}
For fixed $V,s$, every phase test is independent of $n$.  A spatial
factor from output $\alpha$ in slot $i$ is evaluated at
\begin{equation}\label{eq:channel-affine-times}
                T^{e_\alpha i n+a_{\alpha,i}}x_\alpha,
       \qquad e_\alpha=hq_\alpha\in\N,
\end{equation}
where the integer offsets $a_{\alpha,i}$ range over fixed finite sets
as $V,s$ vary.  The original coordinate in slot $b$ is evaluated at
$T^{hb n+a_0}x_0$, again with finitely many possible offsets.

If a selected length is $N=c2^k$, replace it by the union of its
$m=\lfloor N/h\rfloor$ complete intervals of length $h$ and normalize
the resulting sum by $m$.  The error in an average of a unit-bounded
test is at most $2h/N$.  It tends to zero uniformly over all selected
lengths as $a\to\infty$.  We can thus integrate $s$ uniformly over
$[0,h)$ and study averages of $m$ terms.  Discarding any fixed number
of initial indices has the same uniformly vanishing cost.

\emph{Proper collections of groups.}
Under \eqref{eq:channel-proper}, the spatial variables in these groups
are independent of $x_0$ and of one another, conditionally on $V$.
The length $N_a$ is a function of $(x_0,u_0)$, with $u_0$ fixed by $V$.
If at least one output factor is centered, the expectation of the
output product tends to zero as $n\to\infty$: first factor it across
the independent input outputs, then use mixing of all orders for the
distinct slot times belonging to each individual output.  Here and
below we use the main theorem of \cite{MixingCompanion} for fixed
bounded functions.  The gaps
in \eqref{eq:channel-affine-times} diverge within each output, and
there are only finitely many offset patterns.  This convergence is
uniform in $V,s$.  More explicitly, let $R_n(V,s)$ denote that output
expectation and choose $n_0$ so that $|R_n(V,s)|\le\delta$ for
$n\ge n_0$.  Conditional on $x_0,V,s$, the absolute value of the
output-averaged tested sum is at most
\[
 \frac1{m(x_0,V)}\sum_{n<m(x_0,V)}|R_n(V,s)|
       \le\delta+\frac{n_0}{\lfloor2^a/h\rfloor}.
\]
Here the original-coordinate test is bounded by one, and the
denominator is positive for large $a$.  The bound handles both the
selected cutoff and its normalization; integration and then
$a\to\infty$, $\delta\downarrow0$ give zero limit.

If all output factors equal $1$ and the original-coordinate test is
centered, Birkhoff's theorem for the ergodic power $T^{hb}$ gives
zero almost-sure limit of its averages, simultaneously for the
finitely many offsets.  The selected $m$ tends to infinity at every
input point, so bounded convergence proves zero limit also for the
selected averages.  This establishes the channel rule involving
the original coordinate.

\emph{All three groups.}
There is no original-coordinate test in this case.  Terms in which
one group has no centered occurrence were just handled by the
proper-group argument.  Suppose, then, that every group contains a
centered occurrence.  For fixed $V,s$ let $\Omega_j$ be the product
of copies of $(X,\mu)$ for the input outputs in group $j$, and write
the three group products as
\[
 x_n\in L^2(\Omega_1),\qquad
 y_n\in L^2(\Omega_2),\qquad
 w_n\in L^2(\Omega_3).
\]
All have norms at most one.  Before a length event is imposed, the
conditional law of these three input groups given $V$ is their
product law.  Phase weights remain outside these Hilbert-space vectors:
the spatial functions are fixed, and their dependence on $V,s$ is only
through the finitely many integer offsets in
\eqref{eq:channel-affine-times}.

For any fixed $\eta>0$, the families $x_n,y_n$ satisfy the small-pair
hypothesis of Theorem~\ref{thm:packing}, after finitely many initial
indices are removed.  To see this, expand a pair inner product into
expectations over the independent outputs in that group.  One output
has a centered occurrence.  Mixing of all orders makes its expectation
small when the distinct times from the two indices are all sufficiently
far apart.  Within either index, this holds for large $n$ and $n'$.
A remaining obstruction requires
\[
              |e_\alpha i n-e_\alpha i'n'+a|\le R
\]
for one of finitely many choices of $i,i',a$ and a fixed $R$.
For each $n$ there are boundedly many $n'$ satisfying any such
inequality, since $e_\alpha i'>0$.  The obstructing pairs therefore
form a graph of bounded degree $L$.  The constants and the initial
cutoff depend on $\eta$ and the fixed tests; the finite offset
patterns make them uniform in $V,s$.

Condition further on an event
$E_{k,c}=\{N_a=c2^k\}$ with conditional probability
$p_{k,c}=\Prob(E_{k,c}\mid V)>0$.  Repeated values of $c2^k$ are assigned
one index so that these events partition the space.  Write $p=p_{k,c}$.
The event is determined by $x_0$ once $V$ is fixed.  By
\eqref{eq:channel-proper}, conditioning on it preserves the product
marginal on any two groups and the marginal on each single group.
The three-group law need not remain a product.  Its density $\rho$
with respect to the original three-group product measure satisfies
\begin{equation}\label{eq:channel-density}
              0\le p\rho\le1,\qquad \int p\rho=p,
                         \norm{p\rho}_2\le\sqrt p.
\end{equation}
Indeed, for any measurable set of group coordinates its probability
jointly with $E_{k,c}$ is at most its original product probability.

For the $m$ terms corresponding to this length, put
$z_n=x_n\otimes y_n$.  We choose a dimension small enough that the
packing bound is $o(m)$ and large enough that its reciprocal is
summable over the length classes.  The choice
$d_m=\lceil(1+\log m)^{3/2}\rceil$ has both properties: for
$m\asymp2^k$, its reciprocal is $O((1+k)^{-3/2})$.
Let $P_m$ project onto $d_m$ orthonormal
eigenvectors of $\sum_n z_nz_n^*$ corresponding to its largest
eigenvalues, counted with multiplicity.  If the range has smaller dimension,
take the whole range; inside a tied eigenspace make any fixed choice.
Fix $0<\varepsilon\le1$, and choose
$0<\eta\le\varepsilon^2/10^4$ first.
Equation~\eqref{eq:packing-average} gives
\begin{equation}\label{eq:channel-leading-average}
              \frac1m\sum_n\norm{P_mz_n}
                           \le\varepsilon+o_{\varepsilon,L}(1).
\end{equation}
If an initial cutoff of $n_0$ indices was needed, apply the general
bound \eqref{eq:packing-quantitative} to the $m-n_0$ remaining vectors
with the original rank bound $d_m$.  The numerator in that bound is
still $m^{o(1)}$, and the discarded indices cost at most $n_0/m$.
This proves \eqref{eq:channel-leading-average} with the cutoff as
well.  All norms here are for the original product measures.

Under the event law the first two groups still have their product
marginal, and the third retains its original marginal.  Cauchy--Schwarz
under this law therefore bounds the contribution from a term
$(P_mz_n)w_n$, including the probability of the event, by
$p\norm{P_mz_n}$.  Thus its averaged contribution is bounded by
$p$ times the right side of \eqref{eq:channel-leading-average}.

Write $z'_n=(1-P_m)z_n$.  The covariance has trace at most $m$, so its
$(d_m+1)$-st eigenvalue is at most $m/(d_m+1)$ whenever a residual is
present.  The residual Gram matrix $A=(\ip{z'_n}{z'_{n'}})_{n,n'}$
has norm at most $m/d_m$.  The Gram matrix
$B=(\ip{w_n}{w_{n'}})_{n,n'}$ is positive with trace at most $m$.
Consequently
\begin{equation}\label{eq:channel-residual}
 \left\|\frac1m\sum_n z'_n\otimes w_n\right\|_2^2
       =\frac{\tr(AB)}{m^2}\le\frac1{d_m}.
\end{equation}
Combining \eqref{eq:channel-density} and
\eqref{eq:channel-residual} bounds the residual contribution from
this event by $\sqrt{p/d_m}$.

It remains to sum over every possible length.  For the fixed integer
$h$, we have $m=\lfloor c2^k/h\rfloor\asymp2^k$ at large $k$,
uniformly for $c\in F$.  Thus $d_m\asymp(1+k)^{3/2}$, and
\begin{equation}\label{eq:channel-tail}
 \sum_{k\ge a,\ c\in F}\sqrt{p_{k,c}/d_{m(k,c)}}
 \le\left(\sum_{k,c}p_{k,c}\right)^{1/2}
      \left(\sum_{k\ge a,\ c\in F}d_{m(k,c)}^{-1}\right)^{1/2}
                 \longrightarrow0.
\end{equation}
The first factor is one, and the second is a tail of a convergent
series.  The leading-space contributions sum to at most
$\varepsilon+o_{\varepsilon,L}(1)$ because the event probabilities
sum to one.  First let $a\to\infty$, then let $\varepsilon\downarrow0$.
This proves the all-group product rule.

All estimates were uniform over the finitely many offset patterns.
For each pattern and each integer $m$, the finite-dimensional
projections can be chosen once, so their use causes no measurability
problem.  We may now integrate in $V,s$ and restore the phase test.
The uniformly vanishing endpieces do not affect the selected law.
Both defining channel identities have been proved.
\end{proof}

\begin{corollary}\label{cor:channel-array-distal}
In Proposition~\ref{prop:channel-closure}, let
$P=(\omega_Y(y_i))_{i\in I}$ be the array of old distal parameters.
The spatial coordinates of all copied outputs are independent with
their product law conditionally on $P$.  The corresponding assertion
holds after any finite iteration of that proposition, conditionally
on the complete array of old $W_Y$-parameters.
\end{corollary}

\begin{proof}
Lemma~\ref{lem:parameter-distal} places the old parameter array in
the maximal distal factor of the new flow, including after any finite
iteration.  Repeated application of
Proposition~\ref{prop:channel-closure} supplies its channel structure.
Lemma~\ref{lem:channel-distal} gives the conditional product law over
that maximal distal factor, hence over the old parameter array.
\end{proof}

The channel and its tests are fixed before the selected limit is
taken.  The packing estimate and the summation over length classes
require no uniform mixing estimates over different channels or over
all iteration depths.
